
%%%%%%%%%%%%%%%%%%%%%%%%%%%%%%%%%%%%
% Recap/Agenda 
%%%%%%%%%%%%%%%%%%%%%%%%%%%%%%%%%%%%
% TODO better formatting
\begin{frame}
    \frametitle{Module 4: Inference}
    \begin{itemize}
        \item \hl{Previously: }Paired data (Chapter 7.2)
        \item \hl{This time: }Difference of two means (Chapter 7.3)
        \item \hl{Reading: }
        \begin{itemize}
            \item Chapter 7.4 for next time
            \item After that: https://sites.stat.columbia.edu/gelman/research/published/power5r.pdf
        \end{itemize}
        \item \hl{Deadlines/Announcements: }Q3 in discussions this week
    \end{itemize}
    
\end{frame}
  
%%%%%%%%%%%%%%%%%%%%%%%%%%%%%%%%%%%%
% Learning objectives:
%%%%%%%%%%%%%%%%%%%%%%%%%%%%%%%%%%%%
\begin{frame}
    \frametitle{Learning Objectives}
    \begin{itemize}
        \item \textbf{M4, LO5: Explain and Use the t-Distribution:} Explain how the t-distribution differs from the normal distribution and why it is used for population mean inference.
        \item \textbf{M4, LO6: Conduct and Interpret t-Tests:} Design, execute, and interpret t-tests for a single population mean, a difference of paired means, and a difference of independent means, calculating the standard error appropriately for each. Describe how to obtain a p-value for a t-test and a critical t-score for a confidence interval.    
    \end{itemize}
\end{frame}
